\chapter{Metrados, presupuesto y viabilidad}
\section{Herramientas y fronteras del metrado}
El presupuesto se genera desde parámetros canónicos mediante Python, con salida CSV editable. Se evita transcribir totales de capítulos distintos. Las longitudes, equipos, usuarios, auxiliares y potencias son las mismas del modelo de diseño. Se conserva la memoria de cálculo y el manifiesto de fuentes; cualquier modificación de diámetro, grupos o cobertura obliga a reconstruir cantidades, balance y costos.

Para obras civiles se estiman espesores y sobreanchos académicos. El canal produce 101.25 \(\mathrm{m^3}\) de concreto, las dos cámaras de desarenación 143.43 y la cámara de carga 34.28, total aproximado 278.95. El refuerzo se estima en 100 \(\mathrm{kg/m^3}\); no representa planos de armado. La casa de 240 \(\mathrm{m^2}\) se valora por superficie construida, evitando sumarle de nuevo estructura y concreto que ya incluye su partida. Las cimentaciones de silletas y anclajes son provisionales y requieren geotecnia.

La tubería corresponde a 350 m de acero de diámetro exterior 1032 mm y espesor nominal 16 mm. Se agregan 5\% de masa para accesorios de fabricación como estimación. Hay 89 apoyos, cuatro conjuntos principales de válvulas, tres grupos hidráulicos, tres de respaldo, T1 de 800 kVA y auxiliar de 20 kVA. Las redes incluyen 8 km primarios, 9 km secundarios, seis subestaciones, 1037 conexiones potenciales y 226 puntos de alumbrado. Las partidas de redes instaladas incluyen su alcance de montaje y no se recargan como equipos sueltos.

\section{Precios y fecha de referencia}
Los precios unitarios de las partidas son \textbf{estimaciones académicas en soles a octubre de 2026}, no cotizaciones entregadas en Limbani. Se conserva el código, unidad, cantidad, precio, estado y criterio de cada partida en CSV. La consulta minorista pública recuperada permite contrastar cemento Sol tipo I de 42.5 kg a S/30.90 en la página de Sodimac; el producto Yura HS mostrado por el mismo sitio aparece sin existencias. Esa referencia no prueba disponibilidad, transporte ni precio de concreto puesto en obra.

No se recuperaron ofertas verificables para los grupos, tubería fabricada, transporte o subestaciones del proyecto. Se usan valores redondeados de estudio y se marca su condición estimada. Antes de comparar inversión real se solicitarán cotizaciones con impuestos, moneda, fecha, garantía, transporte, montaje y capacidad a la altitud requerida. No se añade un IGV general a valores cuyas bases de impuestos aún no están desagregadas; el presupuesto debe rehacerse homogéneamente con ofertas netas e impuestos para una contratación.

El informe IRENA 2025 sobre costos de generación en 2024 aporta comparadores internacionales: para pequeña hidro, promedios ponderados 2018--2024 de aproximadamente USD 2024/kW 1873 en China, 2534 en Brasil, 3677 en Europa y 4084 en África. No son cotizaciones peruanas ni incluyen necesariamente la distribución y el respaldo de este escenario. No se convierten mecánicamente a soles de 2026 sin tipo de cambio, inflación y equivalencia de alcance. Su finalidad es identificar qué estudios adicionales hacen falta, no validar una precisión aparente.
\begingroup\small\input{generado/presupuesto}\endgroup
\input{generado/resumen_costos}
El transporte y montaje adicional se aplica al 12\% de suministros de tubería, equipos de generación e instalaciones de central; las redes ya se valoran instaladas. Los estudios y supervisión se estiman al 10\% de la base con logística, y la contingencia al 20\% del subtotal resultante. No sustituyen metrados detallados, programación de obra ni costos de acceso. El total de estudio es aproximadamente \textbf{S/\dato{CapexM} millones}. Se presenta una sensibilidad de \(\pm35\%\), S/\dato{CapexLowM}--\dato{CapexHighM} millones; no es un intervalo probabilístico validado.
\begin{figure}[htbp]\centering\includegraphics[width=.88\linewidth]{figuras/costos-r03.png}\caption{Distribución de costos directos estimados. Los recargos se muestran por separado en el resumen.}\end{figure}

\section{Balance económico transparente}
Se evalúa una cobertura híbrida, con demanda 2026 constante durante 20 años y tasa real supuesta del 10\%. No es una predicción de ventas; evita mezclar una tarifa actual no investigada con ingresos futuros ficticios. La energía hidro servida es \dato{EHidro} MWh/año. La energía disponible pero no absorbida no se cuenta como ingreso. Para cubrir el déficit se estima energía diésel en bornes de \dato{DieselE} MWh/año, con eficiencia de entrega supuesta 0.90.

El consumo específico de combustible se adopta en 0.27 l/kWh y el precio en S/5.50/l, ambos \textbf{hipótesis}; no provienen de una cotización local. Se estiman \dato{DieselLitres} mil litros/año y S/\dato{FuelCostM} millones/año de combustible. El consumo real puede crecer a baja carga y debe verificarse con curvas de cada grupo, fletes y almacenamiento. Se propone usar Facilito de OSINERGMIN y ofertas de abastecimiento para sustituir ese precio. La operación y mantenimiento sin combustible se aproxima al 3\% de inversión, S/\dato{OpexM} millones/año; esta simplificación debe reemplazarse por un plan de activos y reposiciones.
\input{generado/viabilidad}
Los ingresos ensayados no son tarifas aprobadas ni recomendaciones de cobro. El valor de equilibrio simplificado es aproximadamente S/\dato{Lcoe}/kWh, incluyendo inversión anualizada, operación y combustible para cubrir el perfil supuesto. No se presenta como un LCOE bancable: faltan reposiciones, estructura financiera, impuestos consistentes, inflación local y riesgos de disponibilidad.

\section{Interpretación y estrategia por etapas}
La rentabilidad financiera de la solución completa es desfavorable en los ingresos considerados. Aumentar potencia instalada no elimina el déficit de estiaje. El tercer grupo hidráulico cubre mejor el pico pero aprovecha poca energía adicional, por lo que 600 kW no queda demostrado como óptimo económico. El padrón y las mediciones pueden justificar menor capacidad y una estrategia gradual.

Se recomienda estudiar primero rehabilitación y reutilización técnicamente verificable de activos, confirmar usuarios permanentes, gestionar demanda y comparar dos grupos de 200 kW con respaldo y una combinación solar--almacenamiento. No se cuantifica ahorro por reutilización ni generación solar sin inspección y recurso correspondiente. Una evaluación social podría considerar beneficios de continuidad, salud y actividad económica, pero requeriría demanda real, criterios públicos aplicables y financiamiento identificado; no se inventa una subvención que vuelva positivo el resultado.
